\item Una tira bimetálica de longitud $L$ está hecha de dos listones de diferentes metales acomodados juntos. 
(a) Primero suponga que la tira originalmente es recta. A medida que la tira se calienta, el metal con el mayor coeficiente de expansión promedio se expande más que el otro, lo que fuerza a la tira en un arco con el radio exterior que tiene mayor circunferencia (figura P6.38). Obtenga una expresión para el ángulo de doblado $\theta$ como función de la longitud inicial de las tiras, sus coeficientes de expansión lineal promedio, el cambio en temperatura y la separación de los centros de las tiras ($\Delta r = r_2 - r_1$). 
(b) Demuestre que el ángulo de doblado disminuye a cero cuando $\Delta T$ tiende a cero y también cuando los dos coeficientes de expansión promedio se vuelven iguales. 
(c) \textbf{¿Qué pasaría si?} ¿Qué ocurre si la tira se enfría?
